% =============================================================
%  04_higher_order.tex
%  Capitolo 4 --- Higher-order e closure
% =============================================================

\chapter{Higher-order e closure}

\begin{intuizione}
Nel Capitolo~3 hai scritto \fn{map}, \fn{filter} e \fn{reduce}
a mano, con ricorsione esplicita.
Questo capitolo ti mostra che Lisp li ha già incorporati ---
e molto di più.
Le funzioni sono valori di prima classe: si passano, si restituiscono,
si compongono. Questa è la programmazione funzionale vera.
\end{intuizione}

\vspace{4pt}

% ================================================================
\section{Funzioni come valori}
% ================================================================

\begin{concetto}{Funzioni di prima classe}
\keyline{}{In Common Lisp una funzione è un valore come qualsiasi altro.}
\keyline{}{Può essere assegnata a una variabile, passata come argomento,\\
restituita come risultato di un'altra funzione.}

\boxsep
\detail{Una \textbf{funzione di ordine superiore} (higher-order function)
è una funzione che riceve funzioni come argomenti o ne restituisce una.
In Lisp questo è naturale: è lo stesso meccanismo con cui si passa
un numero.}
\end{concetto}

\begin{confronto}
\textbf{Confronto con Java/C.}
In Java pre-8 le funzioni non sono valori: per passare ``comportamento''
si usano interfacce o classi anonime. Con Java 8+ si usano le lambda
(\texttt{x -> x*2}), ma rimangono una funzionalità aggiunta a posteriori.
In C i puntatori a funzione esistono, ma la sintassi è scomoda e
non si possono creare funzioni al volo.
In Lisp questo meccanismo esiste dal 1960 ed è al centro del linguaggio.
\end{confronto}

\argomento{La sintassi: \fn{\#'} e \fn{funcall}}

\begin{concetto}{\fn{\#'} e \fn{funcall}}
\keyline{\fn{\#'nome}}{riferimento a una funzione definita (``sharp-quote'').}
\keyline{\fn{(funcall f arg...)}}{chiama la funzione contenuta nella variabile \fn{f}.}
\keyline{\fn{(apply f lista)}}{chiama \fn{f} con gli elementi della lista come argomenti.}
\end{concetto}

\begin{esempio}[\fn{\#'}, \fn{funcall}, \fn{apply}]
\begin{lstlisting}
;; #' prende il riferimento alla funzione
(defun quadrato (x) (* x x))
#'quadrato         ; => #<FUNCTION QUADRATO>  (non la chiama!)

;; funcall: chiama la funzione in una variabile
(defparameter *f* #'quadrato)
(funcall *f* 5)    ; => 25

;; apply: espande una lista come argomenti
(apply #'+ '(1 2 3 4))   ; => 10
(apply #'max '(3 1 4 1 5))   ; => 5

;; Differenza tra funcall e apply:
(funcall #'+ 1 2 3)       ; => 6  (argomenti diretti)
(apply   #'+ '(1 2 3))    ; => 6  (argomenti come lista)
(apply   #'+ 1 2 '(3 4))  ; => 10 (misti: prefisso + lista finale)
\end{lstlisting}
\end{esempio}

\begin{nota}
Perché \fn{\#'} e non solo \fn{'}?
\fn{'foo} cita il \emph{simbolo} \texttt{foo};
\fn{\#'foo} prende la \emph{funzione} associata a \texttt{foo}.
In Common Lisp simboli e funzioni vivono in spazi di nomi separati
(a differenza di Scheme). \fn{\#'} è la notazione corretta quando
vuoi passare una funzione.
\end{nota}

% ================================================================
\section{Le funzioni higher-order della libreria}
% ================================================================

\begin{concetto}{\fn{mapcar} --- il map}
\keyline{}{Applica una funzione a ogni elemento, raccoglie i risultati.}

\begin{sintassi}
(mapcar funzione lista)\\
(mapcar funzione lista1 lista2 ...)  \textit{; piu' liste in parallelo}
\end{sintassi}
\end{concetto}

\begin{esempio}[\fn{mapcar}]
\begin{lstlisting}
(mapcar #'1+      '(1 2 3 4))       ; => (2 3 4 5)
(mapcar #'oddp    '(1 2 3 4))       ; => (T NIL T NIL)
(mapcar #'quadrato '(1 2 3 4))      ; => (1 4 9 16)

;; Con piu' liste: prende elementi in parallelo
(mapcar #'+ '(1 2 3) '(10 20 30))  ; => (11 22 33)
(mapcar #'list '(a b c) '(1 2 3))  ; => ((A 1) (B 2) (C 3))
\end{lstlisting}
\end{esempio}

\seprule

\begin{concetto}{\fn{remove-if} e \fn{remove-if-not} --- il filter}
\keyline{\fn{remove-if}}{rimuove gli elementi per cui il predicato è vero.}
\keyline{\fn{remove-if-not}}{mantiene gli elementi per cui il predicato è vero.}

\boxsep
\detail{\fn{remove-if-not} è il ``filter'' classico.
\fn{remove-if} è il suo complemento (``anti-filter'').}
\end{concetto}

\begin{esempio}[\fn{remove-if} e \fn{remove-if-not}]
\begin{lstlisting}
;; Tieni solo i pari
(remove-if-not #'evenp '(1 2 3 4 5 6))   ; => (2 4 6)

;; Elimina i negativi
(remove-if #'minusp '(3 -1 4 -1 5 -9 2)) ; => (3 4 5 2)

;; Con predicato personalizzato
(remove-if-not #'(lambda (x) (> x 3))
               '(1 2 3 4 5))              ; => (4 5)
\end{lstlisting}
\end{esempio}

\begin{nota}
Common Lisp moderno preferisce \fn{remove} con \fn{:test} o
\fn{:key} per casi semplici. Le varianti \fn{-if} accettano
predicati arbitrari. Vedi anche \fn{find-if}, \fn{count-if},
\fn{every}, \fn{some} --- stessa struttura, risultato diverso.
\end{nota}

\seprule

\begin{concetto}{\fn{reduce} --- il reduce/fold}
\keyline{}{Riduce una lista a un singolo valore applicando\\
una funzione binaria in modo cumulativo.}

\begin{sintassi}
(reduce funzione lista)\\
(reduce funzione lista :initial-value valore-iniziale)
\end{sintassi}

\detail{\fn{:initial-value} serve per le liste vuote e come
``elemento neutro'' dell'operazione (0 per la somma, 1 per il prodotto).}
\end{concetto}

\begin{esempio}[\fn{reduce}]
\begin{lstlisting}
;; Somma
(reduce #'+ '(1 2 3 4))              ; => 10
(reduce #'+ '() :initial-value 0)   ; => 0  (lista vuota)

;; Prodotto
(reduce #'* '(2 3 4))               ; => 24

;; max di una lista
(reduce #'max '(3 1 4 1 5 9 2))     ; => 9

;; Concatena stringhe
(reduce #'(lambda (a b)
            (concatenate 'string a " " b))
        '("ciao" "mondo" "Lisp"))    ; => "ciao mondo Lisp"

;; Da lista di cifre a numero intero: 1534
(reduce #'(lambda (acc d) (+ (* acc 10) d))
        '(1 5 3 4) :initial-value 0)  ; => 1534
\end{lstlisting}
\end{esempio}

\begin{intuizione}
\textbf{\fn{reduce} è il più potente dei tre.}
Qualsiasi \fn{mapcar} o \fn{remove-if} si può esprimere con \fn{reduce}.
La direzione predefinita è da sinistra a destra (\emph{fold left});
con \fn{:from-end t} si procede da destra a sinistra (\emph{fold right}).
\end{intuizione}

\seprule

\begin{concetto}{\fn{mapcan} --- map + flatten}
\keyline{}{Come \fn{mapcar}, ma concatena i risultati invece di\\
raccoglierli in una lista di liste.}

\boxsep
\detail{La funzione deve restituire una lista per ogni elemento.
\fn{mapcan} è l'equivalente del \texttt{flatMap} di Java/JavaScript.}
\end{concetto}

\begin{esempio}[\fn{mapcan}]
\begin{lstlisting}
;; mapcar vs mapcan
(mapcar  #'list '(1 2 3))   ; => ((1) (2) (3))  lista di liste
(mapcan  #'list '(1 2 3))   ; => (1 2 3)         appiattita

;; Raddoppia elementi
(mapcan #'(lambda (x) (list x x)) '(1 2 3))
; => (1 1 2 2 3 3)

;; Espandi ogni elemento nei suoi vicini
(mapcan #'(lambda (x) (list (1- x) x (1+ x)))
        '(10 20))
; => (9 10 11 19 20 21)
\end{lstlisting}
\end{esempio}

\begin{workbook}
\textbf{Esercizi 69--84}
(\fn{mapcar}, \fn{remove-if-not}, \fn{reduce}, \fn{mapcan},
\fn{sort}, \fn{every}/\fn{some}, composizione, ordinamento)\\[4pt]
Ordine suggerito per gli essenziali: \textbf{69, 70, 72, 78, 85}.\\[4pt]
L'esercizio~72 (somma con \fn{reduce}) è banale ma fondamentale:
assicurati di capire il ruolo di \fn{:initial-value}.
L'esercizio~75 (cifre a intero) e 77 (Horner) sono i più eleganti:
vedono \fn{reduce} usato in modo non ovvio.
\end{workbook}

% ================================================================
\section{Lambda: funzioni anonime}
% ================================================================

\begin{concetto}{\fn{lambda} --- funzione senza nome}
\keyline{}{Crea una funzione al volo, senza darle un nome.}

\begin{sintassi}
(lambda (param1 param2 ...) corpo)
\end{sintassi}

\detail{Una lambda è una funzione come \fn{defun}, ma non viene
registrata nel namespace globale. Si usa quando la funzione serve
solo in un posto e non vale la pena nominarla.}
\end{concetto}

\begin{esempio}[Lambda con \fn{mapcar} e \fn{remove-if-not}]
\begin{lstlisting}
;; Senza lambda: serve definire una funzione ausiliaria
(defun doppio (x) (* 2 x))
(mapcar #'doppio '(1 2 3))   ; => (2 4 6)

;; Con lambda: diretto, nessuna definizione separata
(mapcar #'(lambda (x) (* 2 x)) '(1 2 3))   ; => (2 4 6)

;; Anche piu' argomenti
(mapcar #'(lambda (x y) (+ x y))
        '(1 2 3) '(10 20 30))              ; => (11 22 33)

;; Lambda in remove-if-not
(remove-if-not #'(lambda (x) (> x 3))
               '(1 2 3 4 5 6))             ; => (4 5 6)
\end{lstlisting}
\end{esempio}

\begin{nota}
La notazione \fn{\#'(lambda ...)} è equivalente a \fn{(lambda ...)} in
Common Lisp moderno --- la lambda è già una funzione di prima classe.
La forma con \fn{\#'} è più esplicita e preferita per chiarezza.
\end{nota}

% ================================================================
\section{Closure: funzioni con memoria}
% ================================================================

\begin{concetto}{Closure --- cattura dell'ambiente lessicale}
\keyline{}{Una closure è una funzione che \emph{cattura} le variabili\\
dell'ambiente in cui è stata creata.}
\keyline{}{Le variabili catturate sopravvivono alla funzione che le ha create.}

\boxsep
\detail{Quando una lambda viene creata, ``chiude sopra'' le variabili
locali visibili in quel momento. Ogni chiamata successiva alla lambda
vede e può modificare quelle variabili --- anche se la funzione che le
aveva definite è già terminata.}
\end{concetto}

\begin{tcolorbox}[
  enhanced,
  colback=black!2, colframe=black!20,
  arc=3pt, boxrule=0.4pt,
  left=12pt, right=12pt, top=9pt, bottom=9pt,
  before skip=8pt, after skip=10pt
]
\small\color{black!65}

\textbf{Come funziona la cattura.}

\medskip
Considera questa funzione:

\begin{tcolorbox}[
  enhanced, colback=clBg, colframe=gray!40,
  arc=2pt, boxrule=0.4pt,
  left=8pt, right=6pt, top=4pt, bottom=4pt,
  before skip=4pt, after skip=4pt
]
\begin{lstlisting}
(defun make-adder (n)
  (lambda (x) (+ x n)))   ; n e' catturata

(defparameter *add5*  (make-adder 5))
(defparameter *add10* (make-adder 10))

(funcall *add5*  3)   ; => 8
(funcall *add10* 3)   ; => 13
(funcall *add5*  7)   ; => 12
\end{lstlisting}
\end{tcolorbox}

Quando \fn{make-adder} viene chiamata con \fn{n=5}, crea una lambda
che cattura quella specifica istanza di \fn{n}.
Quando \fn{make-adder} viene chiamata di nuovo con \fn{n=10},
crea una \emph{nuova} lambda con una \emph{nuova} istanza di \fn{n}.
Le due closure sono indipendenti.

\medskip
In Java questo si chiamerebbe ``parametro catturato da una lambda''.
In C++ ``capture by value''. In JavaScript ``closure su una variabile''.
In tutti i casi il meccanismo è lo stesso: la funzione porta con sé
un riferimento alle variabili dell'ambiente in cui è nata.

\medskip
\textbf{Closure con stato mutabile.}
Se la variabile catturata viene modificata con \fn{setf}, la modifica
è visibile a tutte le chiamate successive della stessa closure.
Questo permette di creare \emph{oggetti con stato} usando solo funzioni:

\begin{tcolorbox}[
  enhanced, colback=clBg, colframe=gray!40,
  arc=2pt, boxrule=0.4pt,
  left=8pt, right=6pt, top=4pt, bottom=4pt,
  before skip=4pt, after skip=4pt
]
\begin{lstlisting}
(defun make-counter (&optional (start 0) (step 1))
  (let ((val start))            ; val e' catturata
    (lambda ()
      (let ((current val))
        (setf val (+ val step)) ; modifica lo stato interno
        current))))             ; restituisce il vecchio valore

(defparameter *c* (make-counter))
(funcall *c*)   ; => 0
(funcall *c*)   ; => 1
(funcall *c*)   ; => 2

(defparameter *c2* (make-counter 10 5))
(funcall *c2*)  ; => 10
(funcall *c2*)  ; => 15
\end{lstlisting}
\end{tcolorbox}

Ogni chiamata a \fn{make-counter} crea una \emph{nuova} variabile \fn{val}
indipendente. Due contatori creati separatamente non interferiscono.
Questo è il fondamento del concetto di oggetto: stato + comportamento
incapsulati insieme.
\end{tcolorbox}

\begin{workbook}
\textbf{Esercizi 85--88}
(Composizione, applicazione parziale, contatore, accumulatore)\\[4pt]
Questi quattro esercizi sono il cuore del capitolo.
L'esercizio~85 (\fn{compose}) e l'87 (\fn{make-counter}) sono gli essenziali.\\[4pt]
\textbf{Attenzione all'es.~87}: la funzione restituisce una \emph{funzione},
non un valore. Se \fn{(make-counter)} restituisce un numero sei sulla
strada sbagliata.
\end{workbook}

% ================================================================
\section{Composizione e applicazione parziale}
% ================================================================

\argomento{Composizione di funzioni}

\begin{concetto}{Composizione}
\keyline{}{$f \circ g$ è la funzione che prima applica $g$, poi $f$ al risultato.}
\end{concetto}

\begin{esempio}[Composizione manuale e generalizzata]
\begin{lstlisting}
;; Composizione manuale
(defun componi-due (f g)
  (lambda (x) (funcall f (funcall g x))))

(defparameter *inc-poi-doppio*
  (componi-due #'(lambda (x) (* 2 x)) #'1+))

(funcall *inc-poi-doppio* 5)   ; => 12  ((5+1)*2)


;; Composizione generalizzata: n funzioni
(defun compose (&rest funzioni)
  (if (null funzioni)
      #'identity
      (let ((f (car funzioni))
            (resto (apply #'compose (cdr funzioni))))
        (lambda (x) (funcall f (funcall resto x))))))

(defparameter *pipeline*
  (compose #'(lambda (x) (* x x))   ; terzo: quadrato
           #'1+                      ; secondo: incrementa
           #'abs))                   ; primo: valore assoluto

(funcall *pipeline* -3)   ; => 16   abs(-3)=3 -> 3+1=4 -> 4^2=16
\end{lstlisting}
\end{esempio}

\argomento{Applicazione parziale}

\begin{concetto}{Applicazione parziale}
\keyline{}{Fissa alcuni argomenti di una funzione, restituisce\\
una funzione che attende i rimanenti.}

\boxsep
\detail{In linguaggi funzionali puri (Haskell) ogni funzione è
automaticamente ``currificata''. In Common Lisp si costruisce
esplicitamente con una closure.}
\end{concetto}

\begin{esempio}[Applicazione parziale]
\begin{lstlisting}
(defun partial (f &rest args-fissi)
  (lambda (&rest args-nuovi)
    (apply f (append args-fissi args-nuovi))))

;; Costruisce "add-10" fissando il primo argomento di +
(defparameter *add10* (partial #'+ 10))
(funcall *add10* 5)    ; => 15
(funcall *add10* 20)   ; => 30

;; Costruisce un moltiplicatore
(defparameter *doppio* (partial #'* 2))
(mapcar *doppio* '(1 2 3 4))   ; => (2 4 6 8)

;; Combinato con mapcar: filtra i multipli di 3
(defparameter *mult3?* (partial #'= 0))
;; Non funziona direttamente, ma l'idea e' questa:
(remove-if-not #'(lambda (x) (= (mod x 3) 0))
               '(1 2 3 4 5 6 7 8 9))   ; => (3 6 9)
\end{lstlisting}
\end{esempio}

% ================================================================
\section{Funzioni utili della libreria}
% ================================================================

\begin{esempio}[Altre funzioni higher-order della libreria]
\begin{lstlisting}
;; sort: ordina con comparatore custom (DISTRUTTIVO!)
(sort '(3 1 4 1 5 9) #'<)    ; => (1 1 3 4 5 9)
(sort '(3 1 4 1 5 9) #'>)    ; => (9 5 4 3 1 1)

;; sort su strutture con :key
(sort '((anna . 92) (mario . 85) (luca . 78))
      #'>
      :key #'cdr)
; => ((ANNA . 92) (MARIO . 85) (LUCA . 78))

;; every / some: predicati su sequenze
(every  #'evenp '(2 4 6))    ; => T   (tutti pari)
(every  #'evenp '(2 3 6))    ; => NIL
(some   #'oddp  '(2 3 4))    ; => T   (almeno uno dispari)
(notany #'minusp '(1 2 3))   ; => T   (nessun negativo)

;; count-if: conta gli elementi che soddisfano il predicato
(count-if #'evenp '(1 2 3 4 5 6))   ; => 3

;; find-if: trova il primo elemento che soddisfa il predicato
(find-if #'evenp '(1 3 4 5 6))      ; => 4
\end{lstlisting}
\end{esempio}

\begin{attenzione}
\textbf{\fn{sort} è distruttivo.}
Modifica la lista originale. Se hai bisogno della lista originale
intatta, usa \fn{sort} su una copia: \fn{(sort (copy-list lst) \#'<)}.
Questa è una delle poche funzioni della libreria di liste che modifica
la struttura in-place.
\end{attenzione}

% ================================================================
\section*{Riepilogo del capitolo}
% ================================================================

\begin{tcolorbox}[
  enhanced,
  colback=csBlue!5, colframe=csBlue!40,
  arc=4pt, boxrule=0.6pt,
  left=14pt, right=14pt, top=12pt, bottom=12pt,
  before skip=16pt, after skip=16pt
]
\textbf{Quello che hai imparato:}

\medskip
\begin{itemize}
  \item \textbf{Funzioni di prima classe}: \fn{\#'nome} prende il
        riferimento, \fn{funcall} chiama, \fn{apply} espande una lista
        come argomenti.
  \item \textbf{\fn{mapcar}}: trasforma ogni elemento. Con più liste,
        prende elementi in parallelo.
  \item \textbf{\fn{remove-if}/\fn{remove-if-not}}: filtra per predicato.
  \item \textbf{\fn{reduce}}: riduce una lista a un valore.
        \fn{:initial-value} per le liste vuote.
  \item \textbf{\fn{mapcan}}: map + flatten. Equivalente a \texttt{flatMap}.
  \item \textbf{Lambda}: funzione anonima, creata al volo.
  \item \textbf{Closure}: funzione che cattura variabili dell'ambiente.
        Sopravvivono alla funzione che le ha create.
        Con \fn{setf} permettono stato mutabile incapsulato.
  \item \textbf{Composizione e parziale}: costruire nuove funzioni
        combinando quelle esistenti.
  \item \textbf{\fn{sort}, \fn{every}, \fn{some}, \fn{count-if},
        \fn{find-if}}: la libreria standard higher-order.
\end{itemize}

\medskip
\textbf{Prossimo:} Capitolo~5 --- \emph{Strutture dati}.
Association list, hash table, \fn{defstruct}, CLOS base.
\end{tcolorbox}

\begin{workbook}
\textbf{Prima di andare al Capitolo~5:}\\[4pt]
Completa gli esercizi \textbf{69--88} (capitolo \emph{Higher-order
e closure}).\\[4pt]
Priorità: \textbf{69, 70, 72, 78, 85, 87} --- questi sei coprono
tutti i pattern fondamentali del capitolo.\\[4pt]
Se l'esercizio~85 (\fn{compose}) o l'87 (\fn{make-counter}) ti
mette in difficoltà, leggi il suggerimento prima della soluzione.
\end{workbook}
